% ===================== Chapter 19 =====================
\chapter{The Foreign Function Interface}
\label{ch:ffi}

Salam compiles to C, which gives it a superpower: it can call C functions
directly, with no glue code and no performance penalty. The decades of C
libraries the world runs on math, compression, databases, graphics, the
operating system itself are all reachable from Salam through the
\textbf{foreign function interface}, or FFI. This chapter shows how to declare a C
function and call it as if it were your own.

\section{Declaring a C function}

To call a C function you first tell Salam about it: its name, its parameters, and
its return type. You do this in an \code{extern "C"} block, which lists the
\emph{signatures} of C functions without bodies the bodies live in C, after
all:

\begin{salam}
extern "C":
    func printf(format str, ...) i32
    func strlen(s str) u64
end

func main:
    printf("Hello from C's printf!\n")
    msg str = "external"
    printf("the length of \"%s\" is %d\n", msg, strlen(msg))
end
\end{salam}

\begin{progout}
Hello from C's printf!
the length of "external" is 8
\end{progout}

Inside the \code{extern "C"} block, each \code{func} line is a \emph{declaration}:
it says ``a C function with this name and shape exists; let me call it.'' There is
no body. When you call \code{printf} or \code{strlen} from Salam, the compiler
emits a direct call to the real C function exactly as a C program would. The
trailing \code{...} on \code{printf} marks it \textbf{variadic} (it accepts a
variable number of arguments), which is how \code{printf} takes its format
arguments.

\section{How Salam types map to C}

For an FFI call to work, Salam's types must line up with C's. The mapping is
direct and is why the integration is seamless:

\begin{center}
\begin{tabular}{@{}lll@{}}
\toprule
\textbf{Salam} & \textbf{C} & \textbf{Notes} \\
\midrule
\code{i32} & \code{int32\_t} & a 32-bit signed integer \\
\code{i64} & \code{int64\_t} & a 64-bit signed integer \\
\code{u64} & \code{uint64\_t} & a 64-bit unsigned integer \\
\code{f64} & \code{double} & a 64-bit float \\
\code{f32} & \code{float} & a 32-bit float \\
\code{bool} & one byte & \\
\code{char} & one byte & \\
\code{str} & \code{const char*} & a NUL-terminated C string \\
\code{void*} & \code{void*} & an untyped pointer \\
(no return type) & \code{void} & a procedure \\
\bottomrule
\end{tabular}
\end{center}

When you declare an \code{extern} function, choose the Salam types that match the C
function's real signature according to this table. \code{strlen} in C takes a
\code{const char*} and returns a \code{size\_t}, so we declared it as
\code{func strlen(s str) u64}.

\section{Calling the C math library}

A practical example: C's math library (\code{libm}) has fast, accurate
implementations of functions like \code{sqrt}, \code{pow}, and \code{hypot}.
Declare the ones you want and call them:

\begin{salam}
extern "C":
    func printf(format str, ...) i32
    func sqrt(x f64) f64
    func pow(base f64, exp f64) f64
    func hypot(x f64, y f64) f64
end

func main:
    printf("sqrt(2)    = %.6f\n", sqrt(2.0))
    printf("2^10       = %.0f\n", pow(2.0, 10.0))
    printf("hypot(3,4) = %.1f\n", hypot(3.0, 4.0))
end
\end{salam}

\begin{progout}
sqrt(2)    = 1.414214
2^10       = 1024
hypot(3,4) = 5.0
\end{progout}

Each declared function is called like any Salam function, and the result comes
straight back from the C implementation. (Of course, for these particular
functions you would normally just use Salam's own \code{math} module which is
itself partly built on \code{libm} but the same technique works for \emph{any}
C library, including ones with no Salam wrapper.)

\section{Linking external libraries}

The functions above \code{printf}, \code{sqrt} come from the C standard
library, which is always linked in, so declaring them is enough. To use a function
from a \emph{separate} library, you must also tell the compiler to link that
library, with a \code{link} directive before the \code{extern} block:

\begin{salam}
link "sqlite3"

extern "C":
    func sqlite3_libversion() str
    func printf(format str, ...) i32
end

func main:
    printf("SQLite version: %s\n", sqlite3_libversion())
end
\end{salam}

\code{link "sqlite3"} instructs the compiler to link against the SQLite library
(the equivalent of passing \code{-lsqlite3} to a C compiler), making its functions
available to your \code{extern} declarations. You can link several libraries with
several \code{link} lines.

\begin{deep}[In Depth: platform-conditional linking]
Library names sometimes differ across operating systems. Salam's preprocessor lets
you choose the right one at build time with \code{@if}:
\begin{salam}
@if SALAM_OS_WINDOWS
link "winsqlite3"
@else
link "sqlite3"
@end
\end{salam}
The symbols \code{SALAM\_OS\_WINDOWS}, \code{SALAM\_OS\_LINUX}, and
\code{SALAM\_OS\_MACOS} let a single program link the correct library on each
platform the same technique C programmers use with \code{\#ifdef}, expressed in
Salam.
\end{deep}

\section{Working with pointers across the boundary}

Many C functions deal in pointers \code{malloc} returns one, \code{free} takes
one. Salam's \code{void*} type matches C's, so you can call them directly:

\begin{salam}
extern "C":
    func malloc(size u64) void*
    func free(ptr void*)
    func puts(s str) i32
end

func main:
    mut p void* = malloc(64 as u64)
    if p == null:
        puts("malloc failed")
    else:
        puts("malloc(64) gave a non-null pointer")
        free(p)
    end
end
\end{salam}

\begin{progout}
malloc(64) gave a non-null pointer
\end{progout}

Notice that the safe habits from Chapter~18 carry straight over the boundary:
check the pointer against \code{null}, and free what you allocate. The FFI does not
add safety of its own you are calling C, with C's rules so the discipline is
yours to keep.

\section{When to use the FFI}

The FFI is a bridge to be crossed deliberately, not a place to live. Reach for it
when:

\begin{itemize}[leftmargin=1.4em]
  \item you need a capability that exists as a mature C library a database
        driver, a compression codec, an image format and rewriting it in Salam
        would be wasteful;
  \item you must call an operating-system function the standard library does not
        wrap;
  \item you are integrating Salam into an existing C codebase.
\end{itemize}

For everything else, prefer Salam's own standard library and types: they are safer
(the compiler checks them), more convenient, and bilingual. The FFI is the escape
hatch for when you genuinely need the wider C world powerful exactly because it
is there when you need it and ignorable when you do not.

\begin{warn}
Code reached through the FFI is outside Salam's safety net. The compiler trusts
your \code{extern} declarations: if you declare a function with the wrong types,
or pass a bad pointer, the result is a crash or corruption, not a friendly
compile-time error. Declare each C function to match its real signature exactly,
and treat values that cross the boundary with the same care a C programmer would.
\end{warn}

\section{Chapter summary}

\begin{itemize}[leftmargin=1.4em]
  \item An \code{extern "C":} block declares C functions (name, parameters,
        return type, no body) so Salam can call them directly.
  \item Salam types map one-to-one onto C types (\code{i32}$\to$\code{int32\_t},
        \code{str}$\to$\code{const char*}, \code{void*}$\to$\code{void*}, \dots);
        declare each function to match.
  \item A trailing \code{...} marks a variadic function such as \code{printf}.
  \item Standard-library C functions need only declaration; others need a
        \code{link "lib"} directive (optionally guarded by \code{@if} for
        per-platform names).
  \item The FFI adds no safety check pointers, free what you allocate, and
        declare signatures exactly.
  \item Use the FFI to reach mature C libraries and OS functions; prefer Salam's
        own library otherwise.
\end{itemize}

\section{Exercises}

\exercise Declare C's \code{puts} (\code{func puts(s str) i32}) and use it to print
two lines.

\exercise Declare \code{sqrt} and \code{pow} from the math library and print the
square root of 2 and 2 to the 16th power.

\exercise Declare \code{getenv(name str) str} and print the value of an environment
variable of your choice.

\exercise Declare \code{malloc} and \code{free}, allocate 256 bytes, check for
\code{null}, and free them guarding the free with \code{defer}.

\exercise Explain in a comment why an incorrect \code{extern} declaration cannot be
caught by the Salam compiler, and what discipline protects you instead.
